% GENERATED FILE. Edit canonical lessons, then run npm run generate:books.
% canonical-source-hash: fnv1a64:2429d6a5d3742c7b
% canonical-lessons: GU-C539-makkam, GU-C539-tharavvu, GU-C539-kasoti, GU-C539-jhamplavvu, GU-C539-niscay

\chapter{Making Up My Mind}
\label{ch:gu-making-up-my-mind}
\hlchaptermodality{539}

\begin{chapteropening}
\textbf{By the end of this chapter:} \emph{I can say that I have made up my mind.}

Complete the last lesson of chapter 539: I can say that I have made up my mind.
\end{chapteropening}

\section[\texorpdfstring{\gu{મક્કમ} (makkam)}{મક્કમ (makkam)}]{\gu{મક્કમ} (makkam) — resolute}

\label{lesson:GU-C539-makkam}



\begin{quote}
Before the new one: say the Gujarati for to mull over, then the Gujarati for a settlement.
\end{quote}

\subsection*{You'll want to know: \gu{મક્કમ}}
\textbf{\gu{મક્કમ}} — \emph{makkam} — \textquotedblleft{}resolute\textquotedblright{}.

\begin{quote}
\gu{જો} \gu{તમે} \gu{મક્કમ} \gu{રહેશો} \gu{તો} \gu{તમે} \gu{જીતશો}.
\end{quote}

\begin{quote}
\emph{jo tame makkam raheśo to tame jītśo.}
\end{quote}

\begin{quote}
If you stay firm, you will win.
\end{quote}

\begin{cousinweb}
\textbf{\gu{મક્કમ}} (\emph{makkam}) is usually traced to Arabic \emph{muḥkam}, \textquotedblleft{}made firm, strong\textquotedblright{}, through Persian.
\end{cousinweb}

\begin{grammarlens}[title={What you've built}]
One more word for saying what happens if.
\end{grammarlens}

\subsection*{Your turn}
\begin{itemize}
\raggedright
  \item \emph{Say it:} \emph{makkam}
  \item \emph{Say it:} \emph{makkam}, once more
  \item \emph{Say it:} \emph{jo tame makkam raheśo to tame jītśo.} — If you stay firm, you will win.
  \item \emph{Recall:} say the Gujarati for to mull over, then the Gujarati for a settlement, then say \emph{makkam} again
\end{itemize}

\begin{tcolorbox}[breakable,colback=teal!4,colframe=teal!35!black,title={Before you move on}]
What does \gu{મક્કમ} mean? (\textquotedblleft{}Resolute\textquotedblright{}.) Say it once more.
\end{tcolorbox}
\section[\texorpdfstring{\gu{ઠરાવવું} (ṭharāvvũ)}{ઠરાવવું (ṭharāvvũ)}]{\gu{ઠરાવવું} (ṭharāvvũ) — to settle on (a decision)}

\label{lesson:GU-C539-tharavvu}



\begin{quote}
Before the new one: say the Gujarati for a settlement, then the Gujarati for resolute.
\end{quote}

\subsection*{You'll want to know: \gu{ઠરાવવું}}
\textbf{\gu{ઠરાવવું}} — \emph{ṭharāvvũ} — \textquotedblleft{}to settle on (a decision)\textquotedblright{}.

\begin{quote}
\gu{અમે} \gu{ઠરાવ્યું} \gu{છે} \gu{કે} \gu{જો} \gu{વરસાદ} \gu{પડશે} \gu{તો} \gu{અમે} \gu{નહીં} \gu{જઈએ}.
\end{quote}

\begin{quote}
\emph{ame ṭharāvyũ che ke jo varsād paḍśe to ame nahī̃ jaīe.}
\end{quote}

\begin{quote}
We have decided that if it rains we won't go.
\end{quote}

Past, with \gu{તેણે} (he or she …): \gu{તેણે} \gu{ઠરાવ્યું}, \emph{teṇe ṭharāvyũ}.

\begin{cousinweb}
\textbf{\gu{ઠરાવવું}} (\emph{ṭharāvvũ}) is the causative of \textbf{\gu{ઠરવું}}, \textquotedblleft{}to settle, to become fixed\textquotedblright{}, which CDIAL lists under \emph{*stharati}, \textquotedblleft{}is firm\textquotedblright{} (CDIAL 13742): to make a decision firm.
\end{cousinweb}

\begin{grammarlens}[title={What you've built}]
One more word for saying what happens if.
\end{grammarlens}

\subsection*{Your turn}
\begin{itemize}
\raggedright
  \item \emph{Say it:} \emph{ṭharāvvũ}
  \item \emph{Say it:} \emph{ṭharāvvũ}, once more
  \item \emph{Say it:} \emph{ame ṭharāvyũ che ke jo varsād paḍśe to ame nahī̃ jaīe.} — We have decided that if it rains we won't go.
  \item \emph{Recall:} say the Gujarati for a settlement, then the Gujarati for resolute, then say \emph{ṭharāvvũ} again
\end{itemize}

\begin{tcolorbox}[breakable,colback=teal!4,colframe=teal!35!black,title={Before you move on}]
What does \gu{ઠરાવવું} mean? (\textquotedblleft{}To settle on (a decision)\textquotedblright{}.) Say it once more.
\end{tcolorbox}
\section[\texorpdfstring{\gu{કસોટી} (kasoṭī)}{કસોટી (kasoṭī)}]{\gu{કસોટી} (kasoṭī) — a touchstone; a hard test}

\label{lesson:GU-C539-kasoti}



\begin{quote}
Before the new one: say the Gujarati for resolute, then the Gujarati for to settle on.
\end{quote}

\subsection*{You'll want to know: \gu{કસોટી}}
\textbf{\gu{કસોટી}} — \emph{kasoṭī} — \textquotedblleft{}a touchstone; a hard test\textquotedblright{}.

\begin{quote}
\gu{જો} \gu{તમે} \gu{આ} \gu{કસોટીમાં} \gu{મક્કમ} \gu{રહેશો} \gu{તો} \gu{તમને} \gu{આ} \gu{નોકરી} \gu{મળશે}.
\end{quote}

\begin{quote}
\emph{jo tame ā kasoṭīmā̃ makkam raheśo to tamne ā nokrī maḷśe.}
\end{quote}

\begin{quote}
If you stand firm through this test, you will get this job.
\end{quote}

A feminine noun.

\begin{cousinweb}
\textbf{\gu{કસોટી}} (\emph{kasoṭī}) is first the touchstone, the dark stone on which gold is rubbed to show its worth, and so any hard test. CDIAL lists the Gujarati word, \textquotedblleft{}a touchstone\textquotedblright{}, under Sanskrit \emph{kaṣapaṭṭikā} (CDIAL 2973); the sense of a trial it does not give.
\end{cousinweb}

\begin{grammarlens}[title={What you've built}]
One more word for saying what happens if.
\end{grammarlens}

\subsection*{Your turn}
\begin{itemize}
\raggedright
  \item \emph{Say it:} \emph{kasoṭī}
  \item \emph{Say it:} \emph{kasoṭī}, once more
  \item \emph{Say it:} \emph{jo tame ā kasoṭīmā̃ makkam raheśo to tamne ā nokrī maḷśe.} — If you stand firm through this test, you will get this job.
  \item \emph{Recall:} say the Gujarati for resolute, then the Gujarati for to settle on, then say \emph{kasoṭī} again
\end{itemize}

\begin{tcolorbox}[breakable,colback=teal!4,colframe=teal!35!black,title={Before you move on}]
What does \gu{કસોટી} mean? (\textquotedblleft{}A touchstone; a hard test\textquotedblright{}.) Say it once more.
\end{tcolorbox}
\section[\texorpdfstring{\gu{ઝંપલાવવું} (jhamplāvvũ)}{ઝંપલાવવું (jhamplāvvũ)}]{\gu{ઝંપલાવવું} (jhamplāvvũ) — to take the plunge}

\label{lesson:GU-C539-jhamplavvu}



\begin{quote}
Before the new one: say the Gujarati for to settle on, then the Gujarati for a hard test.
\end{quote}

\subsection*{You'll want to know: \gu{ઝંપલાવવું}}
\textbf{\gu{ઝંપલાવવું}} — \emph{jhamplāvvũ} — \textquotedblleft{}to take the plunge\textquotedblright{}.

\begin{quote}
\gu{જો} \gu{હું} \gu{ત્યારે} \gu{ધંધામાં} \gu{ઝંપલાવત} \gu{તો} \gu{આજે} \gu{હું} \gu{અમીર} \gu{હોત}.
\end{quote}

\begin{quote}
\emph{jo hũ tyāre dhandhāmā̃ jhamplāvat to āje hũ amīr hot.}
\end{quote}

\begin{quote}
If I had taken the plunge into business then, I would be rich today.
\end{quote}

Past, with \gu{તેણે} (he or she …): \gu{તેણે} \gu{ધંધામાં} \gu{ઝંપલાવ્યું}, \emph{teṇe dhandhāmā̃ jhamplāvyũ}.

\begin{cousinweb}
\textbf{\gu{ઝંપલાવવું}} (\emph{jhamplāvvũ}) is perhaps related to a reconstructed \emph{*jhapp}, \textquotedblleft{}sudden movement\textquotedblright{}, with a variant \emph{*jhamp-}, \textquotedblleft{}a jump\textquotedblright{}, under which CDIAL lists a Gujarati verb, \textquotedblleft{}to throw headlong\textquotedblright{} (CDIAL 5336); this form is not listed, so its history is not settled. One plunges into water, or into a venture.
\end{cousinweb}

\begin{grammarlens}[title={What you've built}]
One more word for saying what happens if.
\end{grammarlens}

\subsection*{Your turn}
\begin{itemize}
\raggedright
  \item \emph{Say it:} \emph{jhamplāvvũ}
  \item \emph{Say it:} \emph{jhamplāvvũ}, once more
  \item \emph{Say it:} \emph{jo hũ tyāre dhandhāmā̃ jhamplāvat to āje hũ amīr hot.} — If I had taken the plunge into business then, I would be rich today.
  \item \emph{Recall:} say the Gujarati for to settle on, then the Gujarati for a hard test, then say \emph{jhamplāvvũ} again
\end{itemize}

\begin{tcolorbox}[breakable,colback=teal!4,colframe=teal!35!black,title={Before you move on}]
What does \gu{ઝંપલાવવું} mean? (\textquotedblleft{}To take the plunge\textquotedblright{}.) Say it once more.
\end{tcolorbox}
\section[\texorpdfstring{\gu{નિશ્ચય} (niścay)}{નિશ્ચય (niścay)}]{\gu{નિશ્ચય} (niścay) — resolve (a firm decision)}

\label{lesson:GU-C539-niscay}



\begin{quote}
Before the new one: say the Gujarati for a hard test, then the Gujarati for to take the plunge.
\end{quote}

\subsection*{You'll want to know: \gu{નિશ્ચય}}
\textbf{\gu{નિશ્ચય}} — \emph{niścay} — \textquotedblleft{}resolve (a firm decision)\textquotedblright{}.

\begin{quote}
\gu{જો} \gu{તમારો} \gu{નિશ્ચય} \gu{મક્કમ} \gu{હશે} \gu{તો} \gu{તમે} \gu{આ} \gu{કામ} \gu{પતાવી} \gu{દેશો}.
\end{quote}

\begin{quote}
\emph{jo tamāro niścay makkam haśe to tame ā kām patāvī deśo.}
\end{quote}

\begin{quote}
If your resolve is firm, you will finish this work.
\end{quote}

A masculine noun.

\begin{cousinweb}
\textbf{\gu{નિશ્ચય}} (\emph{niścay}) is a tatsama, taken from Sanskrit \emph{niścaya}, \textquotedblleft{}certainty; resolve, a fixed intention\textquotedblright{} (Monier-Williams).
\end{cousinweb}

\begin{grammarlens}[title={What you've built}]
One more word for saying what happens if.
\end{grammarlens}

\subsection*{Your turn}
\begin{itemize}
\raggedright
  \item \emph{Say it:} \emph{niścay}
  \item \emph{Say it:} \emph{niścay}, once more
  \item \emph{Say it:} \emph{jo tamāro niścay makkam haśe to tame ā kām patāvī deśo.} — If your resolve is firm, you will finish this work.
  \item \emph{Recall:} say the Gujarati for a hard test, then the Gujarati for to take the plunge, then say \emph{niścay} again
\end{itemize}

\begin{tcolorbox}[breakable,colback=teal!4,colframe=teal!35!black,title={Before you move on}]
What does \gu{નિશ્ચય} mean? (\textquotedblleft{}Resolve (a firm decision)\textquotedblright{}.) Say it once more.
\end{tcolorbox}
